\todo{Lecture 17}
\section{Extrema liés}
\begin{example}
On considère un cylindre de hauteur $H$ et rayon $R$. On cherche à minimiser la surface, telle que le volume soit égal à $\overline{V}$ pour $R$, $H$. On définit la surface $S(H,R)=2\pi RH+2\pi R^{2}$ et le volume $V(H,R)=\pi R^{2}H=\overline{V}$. On veut $\min_{H,R>0}S(H,R)$ sous la contrainte $\pi R^{2}H-\overline{V}=0$.
\end{example}

Soit $f:E\subset \mathbf{R}^{n}\to \mathbf{R}$, $E$ ouvert non-vide, $f$ est suffisamment régulière (e.g $f\in C^{1}$, voir $C^{2}$). Soit $g:E\to \mathbf{R}$ suffisamment régulière telle que $\Sigma _{g}=\{ \boldsymbol{x}\in E:g(\boldsymbol{x})=0 \}$. On cherche $\min_{\boldsymbol{x}\in E}f(\boldsymbol{x})$ sous la contrainte $g(\boldsymbol{x})=0$ si et seulement si $\min_{\boldsymbol{x}\in\Sigma _{g}}f(\boldsymbol{x})$.

\begin{definition}
On dit que $\boldsymbol{x}^{*}$ est un point de minimum local lié (ou sous contrainte) de $f$ sur $\Sigma _{g}$ si
$$
\exists\delta>0~\forall \boldsymbol{x}\in B(\boldsymbol{x}^{*},\delta)\cap\Sigma _{g}:f(\boldsymbol{x})\ge f(\boldsymbol{x}^{*})
$$
Le minimum est strict si l'inégalité est stricte.

De même pour le maximum local.
\end{definition}
\begin{example}
On cherche $\min_{\boldsymbol{x}\in \mathbf{R}^{2}}(x^{2}+y^{2})$ sous la contrainte $x+y-1=0$. En observant de manière graphique/géométrique le problème, on trouve que le point de minimum est tel que $f$ soit tangentielle à  $\Sigma _{g}$ en ce point. Alternativement, on peut l'interpréter comme le point ou le gradient de $f$ est perpendiculaire à la contrainte.
\end{example}
\begin{example}
On cherche $\min_{\boldsymbol{x}\in \mathbf{R}^{2}}(x+y)$ sous la contrainte $x^{2}+y^{2}-1=0$. Au point de minimum, la courbe de niveau est tangentielle à $\Sigma _{g}$ si et seulement si $\nabla f\perp\Sigma _{g}$ si et seulement si $\nabla f(\boldsymbol{x}^{*})\perp$ plan tangent a $\Sigma _{g}$ en $\boldsymbol{x}^{*}$ si et seulement si $\nabla f(\boldsymbol{x}^{*})\parallel \nabla g(\boldsymbol{x}^{*})$ si et seulement si, il existe $\lambda ^{*}$ tel que $\nabla f(\boldsymbol{x}^{*})=\lambda ^{*}\nabla g(\boldsymbol{x}^{*})$. On peut définir ce plan tangent à $\Sigma _{g}$ comme $\{ \boldsymbol{x}\in \mathbf{R}^{2}:\nabla g(\boldsymbol{x}^{*})(\boldsymbol{x}-\boldsymbol{x^{*}})=0 \}$.
\end{example}
\begin{theorem}[Conditions nécessaires pour extrema locaux lies]
Soit $E\subset \mathbf{R}^{n}$ ouvert non-vide, $f,g:E\to \mathbf{R}$ de classe $C^{1}$. On considère $\Sigma=\{\boldsymbol{z}\in \mathbf{R}^{n}:g(\boldsymbol{z})=0 \}$. Soit $\boldsymbol{z}^{*}\in\Sigma _{g}$ un point d'extremum local lie de $f$ sur $\Sigma _{g}$. Si $\nabla g(\boldsymbol{x}^{*})\neq 0$, alors il existe $\lambda ^{*}\in \mathbf{R}$ tel que $\nabla f(\boldsymbol{z}^{*})=\lambda ^{*}\nabla g(\boldsymbol{z}^{*})$. 
\end{theorem}
\begin{proof}
Puisque $\nabla g(\boldsymbol{z}^{*})\neq0$, il existe $i\in \{ 1,\dots,n \}$ tel que $\frac{ \partial g }{ \partial x_{i} }(\boldsymbol{z}^{*})=0$. Sans perte de généralité, $i=n$, $y=z_{n}$ et $\boldsymbol{x}=(x_{1},\dots,x_{n}-1)=(z_{1},\dots,z_{n-1})\in \mathbf{R}^{n-1}$. Grace au théorème de fonctions implicites, il existe un ouvert $V\subset \mathbf{R}^{n}$ contenant $\boldsymbol{z}^{*}$, un ouvert $U\subset \mathbf{R}^{n-1}$ contenant $\boldsymbol{x}^{*}$ et une fonction unique $\varphi:U\to \mathbf{R}$ de classe $C^{1}$ telle que $\varphi(\boldsymbol{x}^{*})=y^{*}$ et $\mathcal{G}(\varphi)=\Sigma _{g}\cap V$, donc $g(\boldsymbol{x},\varphi(\boldsymbol{x}))=0$ pour tout $\boldsymbol{x}\in U$. On a
$$
\frac{ \partial \varphi }{ \partial x_{i} } (\boldsymbol{x}^{*})=\frac{-\frac{ \partial g }{ \partial x_{i} } (\boldsymbol{z}^{*})}{\frac{ \partial g }{ \partial y } (\boldsymbol{z}^{*})}.
$$
Soit $\tilde{f}(\boldsymbol{x})=f(\boldsymbol{x},\varphi(\boldsymbol{x}))$. La minimisation de cette fonction est équivalente à minimiser $f$ sous la contrainte $g$ car $y=\varphi(\boldsymbol{x})$, donc 
$$
\min_{\boldsymbol{x}\in U}f(\boldsymbol{x},\varphi(\boldsymbol{x}))=\min_{\boldsymbol{x}\in U} \tilde{f}(\boldsymbol{x})
$$
On a que $\boldsymbol{z}^{*}=(\boldsymbol{x}^{*},y^{*})$ est un extremum local de $f$ sur $\Sigma _{g}$ si et seulement si $\boldsymbol{x}^{*}$ est un extremum local de $\tilde{f}$ sur l'ouvert $U$. Alors, nécessairement $\nabla \tilde{f}(\boldsymbol{x}^{*})=0$. $\frac{ \partial f }{ \partial x_{i} }\tilde{f}(\boldsymbol{x}^{*})=0$ pour tout $i\in \{ 1,\dots,n-1 \}$ si et seulement si $$\frac{ \partial f }{ \partial x_{i} }f(\boldsymbol{x}^{*},\varphi(\boldsymbol{x}^{*}))+\frac{ \partial f }{ \partial y }(\boldsymbol{x}^{*},\varphi(\boldsymbol{x}^{*}))\frac{ \partial \varphi }{ \partial x_{i} }(\boldsymbol{x}^{*})=0$$ donc
$$
\frac{ \partial f }{ \partial x_{i} } (\boldsymbol{z}^{*})=\underbrace{ \frac{\frac{ \partial f }{ \partial y } (\boldsymbol{z}^{*})}{\frac{ \partial g }{ \partial y } (\boldsymbol{z}^{*})} }_{ =\lambda ^{*} }  \frac{ \partial g }{ \partial x_{i} } (\boldsymbol{z}^{*})
$$
et
$$
\frac{ \partial f}{ \partial y } (\boldsymbol{z}^{*})=\lambda ^{*}\frac{ \partial g }{ \partial y } (\boldsymbol{z}^{*})
$$
qui conclut la demonstration car $\nabla f(\boldsymbol{z}^{*})=\lambda ^{*}\nabla g(\boldsymbol{z}^{*})$. 
\end{proof}

Pour chercher les extrema lies, on commence par chercher les points $(\boldsymbol{z}^{*},\lambda ^{*})$ tels que $\nabla f(\boldsymbol{z}^{*})=\lambda ^{*}\nabla g(\boldsymbol{z}^{*})$ et $g(\boldsymbol{z}^{*})=0$.

\begin{definition}[Fonction de Lagrange]
On definie le Lagrangien
\begin{align*}
\mathcal{L}:E\times \mathbf{R} & \to \mathbf{R} \\
(\boldsymbol{z},\lambda) & \mapsto f(\boldsymbol{z})-\lambda g(\boldsymbol{z})
\end{align*}
qui a pour gradient
$$
\nabla \mathcal{L}(\boldsymbol{z},^{*})=\begin{pmatrix}
\frac{ \partial \mathcal{L}}{ \partial z_{1} }  \\
\vdots \\
\frac{ \partial \mathcal{L} }{ \partial z_{n} }  \\
\frac{ \partial \mathcal{L} }{ \partial \lambda } 
\end{pmatrix}=\begin{pmatrix}
\frac{ \partial f }{ \partial z_{1} } -\lambda \frac{ \partial g }{ \partial z_{1} } \\
\vdots \\
\frac{ \partial f }{ \partial z_{n} } -\lambda \frac{ \partial g }{ \partial z_{n} } \\
-g(\boldsymbol{z})  
\end{pmatrix}\in \mathbf{R}^{n+1}
$$
On peut resoudre un probleme d'optimisation si et seulement s'il exsite $(\boldsymbol{z}^{*},\lambda ^{*})$ tels que $\nabla \mathcal{L}(\boldsymbol{z}^{*},\lambda ^{*})=\boldsymbol{0}$.
\end{definition}